% =====================================================================
% Sección 6. Adquisición y clasificación de fuentes (Entregable 2)
% Paquetes necesarios en el preámbulo:
%   \usepackage{longtable,booktabs,array,ragged2e,pdflscape,hyperref}
% Reemplaza el Cuadro 1 anterior y el texto que lo introduce.
% =====================================================================

\subsection*{Método de adquisición y criterios de confiabilidad}

El inventario se construyó en tres pasos. Primero, se hizo una búsqueda dirigida en las revistas y el repositorio de la Universidad de Caldas y en bases abiertas (AGRIS, BioOne, Redalyc, Dialnet), con los términos de la ontología: nombre científico de la plaga y del agente, cultivo y técnica analítica. Segundo, se verificó cada fuente en su página de origen para registrar autoría, año, DOI o URL, licencia y fecha de consulta. Tercero, se extrajeron manualmente los datos a una ficha y se asignó cada uno a una clase o propiedad de la ontología.

La confiabilidad usa la escala de la Sección~5.4 (alta, media, baja, insuficiente) y cinco criterios: tipo de publicación, autoría e institución verificables, trazabilidad (DOI o URL estable), pertinencia al dominio (cultivo y región) y accesibilidad. Una fuente puede ser confiable y aun así no aplicar a los cítricos de Caldas; en ese caso se conserva el nivel de la fuente y la limitación se anota de forma explícita. Las fuentes F01 a F17 se detallan en los Cuadros~\ref{tab:inv-meta} y~\ref{tab:inv-conf}; el Cuadro~\ref{tab:inv-cob} contrasta los parámetros del prototipo con las fuentes disponibles.

\makeatletter\g@addto@macro\UrlBreaks{\do\-}\makeatother

\begin{landscape}
\footnotesize
\setlength{\tabcolsep}{4pt}
\renewcommand{\arraystretch}{1.15}
\begin{longtable}{@{}p{0.9cm}p{5.9cm}p{2.6cm}p{2.5cm}p{5.0cm}p{2.3cm}p{1.8cm}@{}}
\caption{Inventario de fuentes: metadatos.}\label{tab:inv-meta}\\
\toprule
\textbf{ID} & \textbf{Referencia} & \textbf{Tipo} & \textbf{Iniciativa} & \textbf{DOI / URL} & \textbf{Acceso y licencia} & \textbf{Consulta}\\
\midrule
\endfirsthead
\toprule
\textbf{ID} & \textbf{Referencia} & \textbf{Tipo} & \textbf{Iniciativa} & \textbf{DOI / URL} & \textbf{Acceso y licencia} & \textbf{Consulta}\\
\midrule
\endhead
\midrule
\multicolumn{7}{r}{\emph{Continúa en la página siguiente}}\\
\endfoot
\bottomrule
\endlastfoot

\RaggedRight F01 & \RaggedRight Arias-Ortega, P. L., Restrepo-García, A. M. y Soto-Giraldo, A. (2016). Primer registro de \textit{Diaphorencyrtus} sp. (Hymenoptera: Encyrtidae) en Colombia. Boletín Científico Centro de Museos, Museo de Historia Natural, 20(1), 157--165. & \RaggedRight Artículo científico (registro) & \RaggedRight Biocontroladores & \RaggedRight \url{https://doi.org/10.17151/bccm.2016.20.1.12} & \RaggedRight Abierto · CC BY 4.0 & 11/nov/2026 \\[3pt]
\RaggedRight F02 & \RaggedRight Restrepo-García, A. M., Arias-Ortega, P. L. y Soto-Giraldo, A. (2016). Primer reporte de \textit{Tamarixia radiata} (Waterston) (Hymenoptera: Eulophidae) en el departamento de Caldas, Colombia. Boletín Científico Centro de Museos, Museo de Historia Natural, 20(1), 204--210. & \RaggedRight Artículo científico (registro) & \RaggedRight Biocontroladores & \RaggedRight \url{https://doi.org/10.17151/bccm.2016.20.1.15} & \RaggedRight Abierto (licencia no verificada) & 11/nov/2026 \\[3pt]
\RaggedRight F03 & \RaggedRight Soto G., A. y Ocampo S., A. (2011). Estudio preliminar de \textit{Trigonospila} sp. (Diptera: Tachinidae), parasitoide de \textit{Compsus viridilineatus} (Coleoptera: Curculionidae). Boletín Científico Centro de Museos, Museo de Historia Natural. & \RaggedRight Artículo científico (descripción) & \RaggedRight Biocontroladores & \RaggedRight \url{https://revistasojs.ucaldas.edu.co/index.php/boletincientifico/article/view/4434} & \RaggedRight Abierto (licencia no verificada) & 11/nov/2026 \\[3pt]
\RaggedRight F04 & \RaggedRight Vinasco, N., Vallejo, L. F. y Soto, A. (2017). New species of the genus \textit{Trigonospila} Pokorny (Diptera: Tachinidae) parasitises adult \textit{Compsus} sp. Schoenherr (Coleoptera: Curculionidae) in Colombia. Boletín Científico Centro de Museos, Museo de Historia Natural, 21(2), 207--219. & \RaggedRight Artículo científico (descripción de especie) & \RaggedRight Biocontroladores & \RaggedRight \url{https://doi.org/10.17151/bccm.2017.21.2.14} & \RaggedRight Abierto (licencia no verificada) & 11/nov/2026 \\[3pt]
\RaggedRight F05 & \RaggedRight Restrepo-García, A. M. y Soto-Giraldo, A. (2017). Control alternativo de \textit{Diaphorina citri} Kuwayama (Hemiptera: Liviidae) utilizando caldo sulfocálcico. Boletín Científico Centro de Museos, Museo de Historia Natural, 21(2), 51--60. & \RaggedRight Artículo científico (laboratorio) & \RaggedRight Biocontroladores & \RaggedRight \url{https://doi.org/10.17151/bccm.2017.21.2.4} & \RaggedRight Abierto (licencia no verificada) & 11/nov/2026 \\[3pt]
\RaggedRight F06 & \RaggedRight Barros-Barrios, M. M., Restrepo-García, A. M., Narváez-Solarte, W. y Soto-Giraldo, A. (2024). Efecto de extractos vegetales compuestos sobre \textit{Diaphorina citri} (Hemiptera: Liviidae) y su parasitoide \textit{Tamarixia radiata} (Hymenoptera: Eulophidae). Boletín Científico Centro de Museos, Museo de Historia Natural, 28(1), 125--138. & \RaggedRight Artículo científico (laboratorio) & \RaggedRight Biocontroladores & \RaggedRight \url{https://doi.org/10.17151/bccm.2024.28.1.7} & \RaggedRight Abierto (licencia no verificada) & 11/nov/2026 \\[3pt]
\RaggedRight F07 & \RaggedRight Restrepo-García, A. M., Pérez-González, O. y Soto-Giraldo, A. (2024). Primer aislamiento de \textit{Hirsutella citriformis} en Colombia y evaluación de su patogenicidad. Boletín Científico Centro de Museos, Museo de Historia Natural, 28(2), 181--190. & \RaggedRight Artículo científico (laboratorio) & \RaggedRight Biocontroladores & \RaggedRight \url{https://doi.org/10.17151/bccm.2024.28.2.9} & \RaggedRight Abierto (licencia no verificada) & 11/nov/2026 \\[3pt]
\RaggedRight F08 & \RaggedRight Gómez-Torres, M. L., Nava, D. E. et al. (2012). Life table of \textit{Tamarixia radiata} (Hymenoptera: Eulophidae) on \textit{Diaphorina citri} (Hemiptera: Psyllidae) at different temperatures. Journal of Economic Entomology, 105(2). & \RaggedRight Artículo científico (laboratorio) & \RaggedRight Biocontroladores & \RaggedRight \url{https://doi.org/10.1603/EC11280} & \RaggedRight Resumen abierto; texto completo posiblemente restringido & 11/nov/2026 \\[3pt]
\RaggedRight F09 & \RaggedRight McCalla, K. A., Keçeci, M., Milosavljevic, I., Ratkowsky, D. A. y Hoddle, M. S. (2019). The influence of temperature variation on life history parameters and thermal performance curves of \textit{Tamarixia radiata} (Hymenoptera: Eulophidae), a parasitoid of the Asian citrus psyllid (Hemiptera: Liviidae). Journal of Economic Entomology, 112(4), 1560--1574. & \RaggedRight Artículo científico (laboratorio) & \RaggedRight Biocontroladores & \RaggedRight \url{https://doi.org/10.1093/jee/toz067} & \RaggedRight Resumen abierto; texto completo posiblemente restringido & 11/nov/2026 \\[3pt]
\RaggedRight F10 & \RaggedRight Monsreal-Ceballos, R. J., Ruíz-Sánchez, E., Sánchez-Borja, M. et al. (2017). Efectos de insecticidas botánicos comerciales en \textit{Tamarixia radiata}, un ectoparasitoide de \textit{Diaphorina citri}. Ecosistemas y Recursos Agropecuarios, 4(12), 589--596. & \RaggedRight Artículo científico (laboratorio) & \RaggedRight Biocontroladores & \RaggedRight \url{https://dialnet.unirioja.es/descarga/articulo/6343377.pdf} & \RaggedRight Abierto (Dialnet; licencia no verificada) & 11/nov/2026 \\[3pt]
\RaggedRight F11 & \RaggedRight Cruz-Cruz, D., Gómez-Ruiz, J., Barrera, J. F., Sánchez-Peña, S. R. y Valle-Mora, J. F. (2020). Pathogenicity of \textit{Beauveria bassiana} against immature stages of \textit{Diaphorina citri} under laboratory conditions. Southwestern Entomologist, 45(3), 609--620. & \RaggedRight Artículo científico (laboratorio) & \RaggedRight Biocontroladores & \RaggedRight \url{https://doi.org/10.3958/059.045.0303} & \RaggedRight Resumen abierto; texto completo por suscripción & 11/nov/2026 \\[3pt]
\RaggedRight F12 & \RaggedRight Imbachi, K., Mesa, N. C., Rodríguez, I. V. et al. (2012). Evaluación de estrategias de control biológico de \textit{Polyphagotarsonemus latus} (Banks) y \textit{Phyllocoptruta oleivora} (Ashmead) en naranja Valencia. Acta Agronómica, 61(4), 333--338. & \RaggedRight Artículo científico (campo) & \RaggedRight Biocontroladores & \RaggedRight \url{https://revistas.unal.edu.co/index.php/acta_agronomica/article/download/38139/40315/169692} & \RaggedRight Abierto (licencia no verificada) & 11/nov/2026 \\[3pt]
\RaggedRight F13 & \RaggedRight Autor por verificar en la ficha del repositorio (2025). Efecto del tratamiento magnético de semillas sobre un sistema productivo de \textit{Solanum lycopersicum} L., a escala comercial. Tesis de Doctorado en Ciencias Agrarias, Universidad de Caldas. & \RaggedRight Tesis de doctorado & \RaggedRight MAGNETIME & \RaggedRight \url{https://repositorio.ucaldas.edu.co/handle/ucaldas/22551} & \RaggedRight Abierto · CC BY-NC-ND 4.0 & 11/nov/2026 \\[3pt]
\RaggedRight F14 & \RaggedRight Torres-Osorio, J. I., Aranzazu-Osorio, J. E. y Carbonell-Padrino, M. V. (2015). Efecto del campo magnético estático homogéneo en la germinación y absorción de agua en semillas de soja. Tecno Lógicas, 18(35), 11--20. & \RaggedRight Artículo científico (laboratorio) & \RaggedRight MAGNETIME & \RaggedRight \url{https://www.redalyc.org/pdf/3442/344243102002.pdf} & \RaggedRight Abierto · CC BY & 11/nov/2026 \\[3pt]
\RaggedRight F15 & \RaggedRight Villa Carmona, E., Torres Osorio, J. I. y Zamorano-Montañez, C. (2023). Análisis del desarrollo vegetativo de plantas de \textit{Solanum lycopersicum} L., generadas de semillas tratadas magnéticamente. Trabajo de grado, Ingeniería Agronómica, Universidad de Caldas. & \RaggedRight Tesis de pregrado & \RaggedRight MAGNETIME & \RaggedRight \url{https://repositorio.ucaldas.edu.co/handle/ucaldas/19679} & \RaggedRight Abierto (openAccess) & 11/nov/2026 \\[3pt]
\RaggedRight F16 & \RaggedRight Universidad de Caldas (2022). Spin Off ClayOMICs de la Universidad de Caldas obtiene su registro de marca. Portal institucional. & \RaggedRight Noticia institucional & \RaggedRight ClayOMICS & \RaggedRight \url{https://www.ucaldas.edu.co/portal/spin-off-clayomics-de-la-universidad-de-caldas-obtiene-su-registro-de-marca/} & \RaggedRight Abierto (web institucional) & 11/nov/2026 \\[3pt]
\RaggedRight F17 & \RaggedRight Betancourt Arango, J. P., Patiño Ospina, A., Taborda Ocampo, G. y Fiscal Ladino, J. A. (2020). Aplicaciones de la xenometabolómica para la identificación de biomarcadores de toxicidad: una revisión del tema. Biosalud, 19(1), 7--30. & \RaggedRight Artículo de revisión & \RaggedRight ClayOMICS & \RaggedRight \url{https://doi.org/10.17151/biosa.2020.19.1.1} & \RaggedRight Abierto (licencia no verificada) & 11/nov/2026 \\[3pt]
\end{longtable}

\begin{longtable}{@{}p{0.9cm}p{6.6cm}p{3.6cm}p{1.4cm}p{3.0cm}p{6.2cm}@{}}
\caption{Inventario de fuentes: extracción y clasificación de confiabilidad.}\label{tab:inv-conf}\\
\toprule
\textbf{ID} & \textbf{Información extraída} & \textbf{Método de extracción} & \textbf{Nivel} & \textbf{Naturaleza} & \textbf{Alcance y limitaciones}\\
\midrule
\endfirsthead
\toprule
\textbf{ID} & \textbf{Información extraída} & \textbf{Método de extracción} & \textbf{Nivel} & \textbf{Naturaleza} & \textbf{Alcance y limitaciones}\\
\midrule
\endhead
\midrule
\multicolumn{6}{r}{\emph{Continúa en la página siguiente}}\\
\endfoot
\bottomrule
\endlastfoot

\RaggedRight F01 & \RaggedRight Primer registro en Colombia de \textit{Diaphorencyrtus} sp., parasitoide de ninfas de \textit{Diaphorina citri}; ejemplares capturados a mano en brotes de cítricos, Manizales. & \RaggedRight Lectura de resumen, métodos y resultados; extracción manual a ficha. Destino: Plaga, AgenteBiocontrol. & \RaggedRight Alta & \RaggedRight Hecho (presencia); hipótesis (control) & \RaggedRight Registro de presencia. No mide parasitismo ni condiciones de uso; el efecto sobre la plaga queda como hipótesis. \\[3pt]
\RaggedRight F02 & \RaggedRight \textit{T. radiata} presente en Caldas, parasitando ninfas de \textit{D. citri} en brotes de cítricos colectados entre enero y abril de 2016 en varios municipios. & \RaggedRight Lectura de resumen, métodos y resultados; extracción manual a ficha. Destino: RelacionBiocontrol (\textit{D. citri}--\textit{T. radiata}), FuenteEvidencia. & \RaggedRight Alta & \RaggedRight Hecho documentado & \RaggedRight Documenta presencia y parasitismo en Caldas; no reporta rangos agroclimáticos ni dosis de liberación. \\[3pt]
\RaggedRight F03 & \RaggedRight Descripción morfológica de larvas, pupas y adultos de \textit{Trigonospila} sp., nuevo parasitoide de adultos del picudo de los cítricos; adultos de \textit{Compsus} recolectados en Maltería, Manizales (2250 msnm). & \RaggedRight Lectura de resumen, métodos y resultados; extracción manual a ficha. Destino: Plaga (\textit{Compsus}), AgenteBiocontrol. & \RaggedRight Alta & \RaggedRight Hecho documentado & \RaggedRight Estudio preliminar morfológico; sin datos de eficacia en campo. Volumen y páginas por completar al citar. \\[3pt]
\RaggedRight F04 & \RaggedRight Descripción de \textit{Trigonospila unicaldasi} sp. nov.; holotipo en Granja Tesorito, Valle de Maltería, Manizales (2280 m), sobre adultos de \textit{Compsus} sp. en limón. & \RaggedRight Lectura de resumen, métodos y resultados; extracción manual a ficha. Destino: AgenteBiocontrol, Plaga, FuenteEvidencia. & \RaggedRight Alta & \RaggedRight Hecho documentado & \RaggedRight Descripción taxonómica y registro de parasitismo; no ofrece condiciones de uso ni protocolo de liberación. \\[3pt]
\RaggedRight F05 & \RaggedRight Toxicidad letal y subletal del caldo sulfocálcico sobre \textit{D. citri} en laboratorio (Centro de Investigación y Cría de Enemigos Naturales, Universidad de Caldas). & \RaggedRight Lectura de resumen, métodos y resultados; extracción manual a ficha. Destino: CondicionUso, restricciones (tieneRestriccion). & \RaggedRight Alta & \RaggedRight Hecho documentado (laboratorio) & \RaggedRight Ensayo de laboratorio sobre la plaga. Verificar si evalúa efecto sobre enemigos naturales antes de usarlo en reglas de incompatibilidad. \\[3pt]
\RaggedRight F06 & \RaggedRight Efecto de extractos de \textit{Murraya paniculata}, \textit{Cocos nucifera}, \textit{Syzygium aromaticum} y \textit{Psidium guajava} sobre \textit{D. citri} y \textit{T. radiata} (olfatometría). & \RaggedRight Lectura de resumen, métodos y resultados; extracción manual a ficha. Destino: tieneRestriccion (compatibilidad con \textit{T. radiata}), Experimento. & \RaggedRight Alta & \RaggedRight Hecho documentado (laboratorio) & \RaggedRight Condiciones controladas. Los valores cuantitativos por extracto deben extraerse del texto completo. \\[3pt]
\RaggedRight F07 & \RaggedRight Cepa colombiana de \textit{H. citriformis} (Granja Montelindo, Universidad de Caldas): mortalidad acumulada de \textit{D. citri} de 98 \% con conidios frente a 3 \% en el control, en bioensayo de laboratorio. & \RaggedRight Lectura de resumen, métodos y resultados; extracción manual a ficha. Destino: AgenteBiocontrol (entomopatógeno), RelacionBiocontrol. & \RaggedRight Alta & \RaggedRight Hecho documentado (laboratorio) & \RaggedRight El artículo reporta incubación a «25 $\pm$ 20 °C» (probable error tipográfico) y no reporta humedad relativa: no sirve para derivar un rango agroclimático. Sin ensayo de campo. \\[3pt]
\RaggedRight F08 & \RaggedRight Tabla de vida de \textit{T. radiata} a 15, 20, 25, 30 y 35 °C (70 $\pm$ 10 \% HR): desarrollo en todas las temperaturas; mayor viabilidad y mayor tasa reproductiva neta a 25 °C. & \RaggedRight Lectura del resumen; extracción manual de rangos de temperatura y humedad. Destino: CondicionUso. & \RaggedRight Alta & \RaggedRight Hecho documentado (laboratorio) & \RaggedRight Población brasileña en laboratorio. «Viabilidad» se refiere al desarrollo, no a la eficacia de una liberación en campo. Completar páginas al citar. \\[3pt]
\RaggedRight F09 & \RaggedRight Efecto de temperaturas constantes y fluctuantes (promedios de 12 a 38 °C) sobre desarrollo y longevidad de \textit{T. radiata}; las fluctuaciones afectan ambos parámetros. & \RaggedRight Lectura del resumen; extracción manual de rangos térmicos. Destino: CondicionUso. & \RaggedRight Alta & \RaggedRight Hecho documentado (laboratorio) & \RaggedRight Población de Pakistán liberada en California; condiciones de laboratorio. Complementa a F08 con un rango térmico más amplio. \\[3pt]
\RaggedRight F10 & \RaggedRight Mortalidad de \textit{T. radiata} por residualidad de 12,7 \% (Biodie) y 12,3 \% (NeemPHC), y por ingestión de 33,3 \% y 28,7 \%; bioensayo a 25 $\pm$ 5 °C y 60 $\pm$ 10 \% HR. & \RaggedRight Lectura de resumen, métodos y resultados; extracción manual a ficha. Destino: tieneRestriccion, TratamientoQuimicoPrevio. & \RaggedRight Alta & \RaggedRight Hecho documentado (laboratorio) & \RaggedRight Evalúa insecticidas botánicos, no organofosforados como el clorpirifos. Estudio hecho en México. \\[3pt]
\RaggedRight F11 & \RaggedRight Cepas de \textit{B. bassiana} (1$\times10^{8}$ conidios/mL): sin efecto sobre huevos; mortalidad de ninfas entre 21 y 58 \% según la cepa; CL50 en ninfas IV--V de 8,34$\times10^{6}$ (Bb-Dc) y 1$\times10^{7}$ (Bb-Rhy) conidios/mL. & \RaggedRight Lectura del resumen; extracción manual de mortalidad y concentraciones. Destino: AgenteBiocontrol, RelacionBiocontrol. & \RaggedRight Alta & \RaggedRight Hecho documentado (laboratorio) & \RaggedRight Cepas mexicanas en laboratorio; no valida la cepa «UC10» del borrador ni aplica directamente a Caldas. \\[3pt]
\RaggedRight F12 & \RaggedRight Ensayo en naranja Valencia (Caicedonia, Valle del Cauca, 1150 msnm): la liberación de \textit{Chrysoperla carnea} redujo el daño de \textit{P. oleivora} a 10 \% en el primer ciclo; los ácaros Phytoseiidae se dispersaron rápido de los frutos. & \RaggedRight Lectura de resumen, métodos y resultados; extracción manual a ficha. Destino: AgenteBiocontrol (\textit{C. carnea}), RelacionBiocontrol. & \RaggedRight Alta & \RaggedRight Hecho documentado (campo) & \RaggedRight Las plagas objetivo son ácaros, no \textit{D. citri}: no sustenta que \textit{C. carnea} controle \textit{D. citri}. Región distinta a Caldas. \\[3pt]
\RaggedRight F13 & \RaggedRight Tratamiento magnético de semillas de tomate a escala comercial: revisión sistemática (PRISMA), optimización bayesiana de parámetros, crecimiento, floración y rendimiento; el rendimiento aumentó y la floración no cambió. & \RaggedRight Lectura de resumen y metadatos del repositorio; extracción manual a ficha. Destino: Experimento, FuenteEvidencia. & \RaggedRight Media & \RaggedRight Hecho documentado (tesis) & \RaggedRight Tesis, no artículo revisado por pares. Cultivo: tomate; no hay validación para cítricos. Corresponde a la «evidencia faltante» que pide el caso integrador. \\[3pt]
\RaggedRight F14 & \RaggedRight Soja expuesta a 125 y 300 mT durante 10, 20 y 60 min; el electroimán superó a los imanes; 300 mT por 10 min redujo más el tiempo medio de germinación; la absorción de agua no mostró diferencias significativas. & \RaggedRight Lectura de resumen, métodos y resultados; extracción manual a ficha. Destino: Experimento (parámetros y resultados). & \RaggedRight Alta & \RaggedRight Hecho documentado (laboratorio) & \RaggedRight Especie: soja. No hay evidencia en cítricos ni en control de plagas. \\[3pt]
\RaggedRight F15 & \RaggedRight Semillas de tomate tratadas con campo magnético estático no homogéneo: plantas con 21 \% más área foliar y 9 \% más altura que el control. & \RaggedRight Lectura de resumen y metadatos del repositorio; extracción manual a ficha. Destino: Experimento. & \RaggedRight Media & \RaggedRight Hecho documentado (tesis) & \RaggedRight Trabajo de grado, sin revisión por pares. Cultivo: tomate. \\[3pt]
\RaggedRight F16 & \RaggedRight ClayOMICS ofrece análisis de datos químicos mediante técnicas de medición, separación e identificación; surge del grupo de investigación en Cromatografía y Técnicas Afines; obtuvo registro de marca ante la SIC en 2022. & \RaggedRight Lectura de la nota; extracción manual de alcance y servicios. Destino: contexto de AnalisisQuimico. & \RaggedRight Media & \RaggedRight Contexto institucional & \RaggedRight Sirve para la existencia y el alcance de los servicios, no como evidencia científica. Fecha exacta de publicación por completar. \\[3pt]
\RaggedRight F17 & \RaggedRight Revisión de técnicas analíticas (GC-MS, LC-MS, UHPLC-QTOF-MS, RMN, ICP-MS) aplicadas a plaguicidas, metales pesados, micotoxinas y otros analitos, en matrices de agua, tejidos, plantas y alimentos. & \RaggedRight Lectura de resumen, métodos y resultados; extracción manual a ficha. Destino: AnalisisQuimico (técnica--analito--matriz). & \RaggedRight Alta & \RaggedRight Hecho documentado (revisión) & \RaggedRight Revisión del grupo de Cromatografía de la Universidad de Caldas, no un servicio ni un resultado de ClayOMICS. Sirve de puente conceptual con residuos de plaguicidas. \\[3pt]
\end{longtable}

\begin{longtable}{@{}p{6.2cm}p{3.4cm}p{2.0cm}p{11.0cm}@{}}
\caption{Cobertura de los parámetros del prototipo por las fuentes del inventario.}\label{tab:inv-cob}\\
\toprule
\textbf{Dato o afirmación del prototipo} & \textbf{Fuentes} & \textbf{Estado} & \textbf{Acción recomendada}\\
\midrule
\endfirsthead
\toprule
\textbf{Dato o afirmación del prototipo} & \textbf{Fuentes} & \textbf{Estado} & \textbf{Acción recomendada}\\
\midrule
\endhead
\bottomrule
\endlastfoot

\RaggedRight \textit{T. radiata} controla a \textit{D. citri} & \RaggedRight F02, F08, F09 & \RaggedRight Documentado & \RaggedRight Mantener la relación con nivel de evidencia alto. \\[3pt]
\RaggedRight Rango térmico 20--32 °C (borrador) o 20--30 °C (PDF) & \RaggedRight F08, F09 & \RaggedRight Parcial & \RaggedRight Las fuentes muestran desarrollo entre 15 y 35 °C con mejor viabilidad a 25 °C (laboratorio). Ningún estudio fija un umbral de «liberación viable»: ajustar el rango a lo documentado o rotularlo sintético. \\[3pt]
\RaggedRight Humedad relativa 60--85 \% & \RaggedRight F08, F10 & \RaggedRight Sin fuente & \RaggedRight Las fuentes usan 70 $\pm$ 10 \% y 60 $\pm$ 10 \% como condiciones de ensayo, no como rango de viabilidad. Rotular como sintético. \\[3pt]
\RaggedRight Dosis de 400--600 individuos por hectárea & \RaggedRight --- & \RaggedRight Sin fuente & \RaggedRight No hay fuente verificada. Rotular como sintético o buscar el protocolo de liberación. \\[3pt]
\RaggedRight Carencia de 21 días tras clorpirifos & \RaggedRight F10 (solo insecticidas botánicos) & \RaggedRight Sin fuente & \RaggedRight F10 documenta mortalidad por insecticidas botánicos, no por organofosforados. Rotular como sintético. \\[3pt]
\RaggedRight Cepa «\textit{Beauveria bassiana} UC10» validada por la Universidad de Caldas & \RaggedRight F11, F07 & \RaggedRight Sin fuente & \RaggedRight F11 usa cepas mexicanas. La cepa colombiana documentada es \textit{Hirsutella citriformis} (F07): usarla o rotular UC10 como sintética. \\[3pt]
\RaggedRight \textit{Chrysoperla carnea} controla a \textit{D. citri} & \RaggedRight F12 & \RaggedRight Sin fuente & \RaggedRight F12 evalúa ácaros en naranja Valencia, no \textit{D. citri}. \\[3pt]
\RaggedRight Ejemplos del PDF: «Estudio Soto 2019», «Gómez 2021», protocolo v2.1, \textit{Aphis gossypii} & \RaggedRight --- & \RaggedRight Sin fuente & \RaggedRight No están en el inventario. Rotular el ejemplo como sintético o retirar las citas. \\[3pt]
\end{longtable}
\end{landscape}
